\documentclass[11pt]{beamer}
\usetheme{Madrid}
\usepackage[T1]{fontenc}
\usepackage[utf8]{inputenc}
\usepackage{amsmath,amsfonts,amssymb}
\usepackage{graphicx}
\usepackage{xcolor}
\usepackage{tikz}
\usepackage{array}
\usepackage{longtable}
\usepackage{booktabs}
\usepackage{hyperref}

\title{EB2-NIW \& EB1A Profile Building}
\subtitle{Research Strategy • Publications • Citations • Peer Review}
\author{Complete DIY Research Profile System}
\date{\today}

\setbeamercovered{transparent}
\setbeamertemplate{navigation symbols}{}
\setbeamerfont{itemize/enumerate body}{size=\scriptsize}
\setbeamerfont{itemize/enumerate subbody}{size=\tiny}
\setbeamerfont{block body}{size=\scriptsize}
\setbeamerfont{block title}{size=\normalsize}

\begin{document}
	
	% ============================================================
	% SLIDE 1: TITLE
	% ============================================================
	\begin{frame}
		\titlepage
		\centering
		\scriptsize \textit{Build Your Research Profile from Zero to EB1A-Ready}
	\end{frame}
	
	% ============================================================
	% SLIDE 2: WHAT IS PROFILE BUILDING?
	% ============================================================
	\begin{frame}{What Is Research Profile Building?}
		\scriptsize
		\begin{block}{Definition}
			Profile building is the strategic process of creating, documenting, and showcasing your research achievements to establish yourself as an expert in your field.
		\end{block}
		\begin{block}{Core Components of a Research Profile}
			\begin{itemize}
				\item \textbf{Publications:} Journal papers, conference proceedings, book chapters
				\item \textbf{Citations:} How many times others reference your work
				\item \textbf{Peer Review:} Judging others' work as a recognized expert
				\item \textbf{Awards \& Recognition:} Prizes, fellowships, honors
				\item \textbf{Indexing:} Your work appearing in major databases (.gov, .edu, Scopus)
				\item \textbf{Media \& Outreach:} News articles, podcasts, public engagement
			\end{itemize}
		\end{block}
	\end{frame}
	
	% ============================================================
	% SLIDE 3: WHY PROFILE BUILDING MATTERS
	% ============================================================
	\begin{frame}{Why Profile Building Matters}
		\scriptsize
		\begin{block}{A Strong Research Profile Demonstrates:}
			\begin{itemize}
				\item You are recognized as an expert by peers
				\item Your work has impact beyond yourself
				\item You contribute to advancing your field
				\item You have sustained achievement over time
			\end{itemize}
		\end{block}
		\begin{block}{What a Strong Profile Enables:}
			\begin{itemize}
				\item Recognition as a leader in your research area
				\item Invitations to collaborate with top institutions
				\item Opportunities to review for prestigious journals
				\item Citations from government agencies and universities
			\end{itemize}
		\end{block}
	\end{frame}
	
	% ============================================================
	% SLIDE 4: PUBLICATION STRATEGY - OVERVIEW
	% ============================================================
	\begin{frame}{Publication Strategy - Complete Overview}
		\scriptsize
		\begin{block}{Three Publication Tiers}
			\centering
			\begin{tabular}{|p{3cm}|p{4cm}|p{3cm}|}
				\hline
				\textbf{Tier} & \textbf{Characteristics} & \textbf{Purpose} \\
				\hline
				Level 0 & New journals, lower fees, fast & Build confidence, start portfolio \\
				\hline
				Mid-Tier & Scopus indexed, moderate fees & Citation building, main volume \\
				\hline
				High-Tier & Q1 Scopus, prestigious journals & Differentiation, top evidence \\
				\hline
			\end{tabular}
		\end{block}
		\begin{block}{Publication Mix Strategy}
			\begin{itemize}
				\item Years 1-2: Focus on Level 0 and Mid-Tier (build volume)
				\item Years 2-3: Add High-Tier papers (demonstrate selectivity)
				\item Ideal mix: 60\% Mid-Tier, 20\% Level 0, 20\% High-Tier
			\end{itemize}
		\end{block}
	\end{frame}
	
	% ============================================================
	% SLIDE 5: LEVEL 0 JOURNALS - GETTING STARTED
	% ============================================================
	\begin{frame}{Level 0 Journals - Your Starting Point}
		\scriptsize
		\begin{block}{What Are Level 0 Journals?}
			\begin{itemize}
				\item Entry-level peer-reviewed journals
				\item May not yet be in Scopus (but working toward it)
				\item Legitimate peer review process
				\item Lower article processing charges (APC \$100-\$500)
				\item Faster publication (4-8 weeks)
			\end{itemize}
		\end{block}
		\begin{block}{Examples by Field (Verify Current Status)}
			\begin{itemize}
				\item \textbf{Computer Science:} International Journal of Advanced Computer Science and Applications
				\item \textbf{Engineering:} Journal of Engineering Science and Technology
				\item \textbf{Multidisciplinary:} International Journal of Recent Technology and Engineering
				\item \textbf{Education:} Journal of Education and Social Research
				\item \textbf{Business:} International Journal of Business and Management
			\end{itemize}
			\textcolor{blue}{Always verify current indexing status before submitting.}
		\end{block}
	\end{frame}
	
	% ============================================================
	% SLIDE 6: MID-TIER JOURNALS - BULK OF YOUR PORTFOLIO
	% ============================================================
	\begin{frame}{Mid-Tier Journals - Building Your Portfolio}
		\scriptsize
		\begin{block}{Characteristics of Strong Mid-Tier Journals}
			\begin{itemize}
				\item Indexed in Scopus (Q2 or Q3)
				\item Moderate APC (\$800-\$1,500)
				\item Published by reputable publishers (MDPI, Springer, Elsevier, Sage, Taylor \& Francis)
				\item 6-12 week peer review
				\item Editorial board includes US researchers
			\end{itemize}
		\end{block}
		\begin{block}{Recommended Mid-Tier Publishers}
			\begin{itemize}
				\item \textbf{MDPI:} Sustainability, Big Data and Cognitive Computing, Future Internet, Applied Sciences
				\item \textbf{Springer:} SN Computer Science, Discover Artificial Intelligence, Discover Data
				\item \textbf{Elsevier:} Array, Software Impacts, Data in Brief
				\item \textbf{Sage:} Sage Open, Journal of Applied Behavioral Science
				\item \textbf{World Scientific:} International Journal of Innovation and Technology Management
			\end{itemize}
		\end{block}
	\end{frame}
	
	% ============================================================
	% SLIDE 7: HIGH-TIER JOURNALS - PRESTIGE EVIDENCE
	% ============================================================
	\begin{frame}{High-Tier Journals - Prestige Building}
		\scriptsize
		\begin{block}{Characteristics of High-Tier Journals}
			\begin{itemize}
				\item Scopus Q1 or Web of Science top quartile
				\item High impact factor (3.0+ depending on field)
				\item Low acceptance rate (10-20\%)
				\item Rigorous peer review (3-6 months)
				\item No or low APC for subscription options
			\end{itemize}
		\end{block}
		\begin{block}{Examples of High-Tier Journals by Field}
			\begin{itemize}
				\item \textbf{Computer Science/AI:} IEEE Transactions on Pattern Analysis and Machine Intelligence, Artificial Intelligence (Elsevier)
				\item \textbf{Finance/Economics:} Journal of Finance, Journal of Financial Economics, Review of Financial Studies
				\item \textbf{General Science:} Nature, Science, PNAS
				\item \textbf{Education:} Computers \& Education, Educational Research Review
				\item \textbf{Cybersecurity:} IEEE Transactions on Information Forensics and Security, Computers \& Security
			\end{itemize}
		\end{block}
	\end{frame}
	
	% ============================================================
	% SLIDE 8: CONFERENCE PROCEEDINGS - FASTER PATHWAY
	% ============================================================
	\begin{frame}{Conference Proceedings - Faster Publications}
		\scriptsize
		\begin{block}{Why Conferences Matter}
			\begin{itemize}
				\item Faster publication (3-6 months vs 12+ months for journals)
				\item Lower cost (registration \$300-\$800 vs APC \$1,000+)
				\item Networking opportunities with researchers
				\item Presentation experience (public speaking)
			\end{itemize}
		\end{block}
		\begin{block}{Prestigious Conference Proceedings}
			\begin{itemize}
				\item \textbf{IEEE Xplore:} All IEEE conference proceedings (Scopus indexed)
				\item \textbf{Springer LNCS/LNAI/CCIS:} Conference proceedings published by Springer
				\item \textbf{ACM Digital Library:} ACM conference proceedings
				\item \textbf{Elsevier Procedia:} Procedia Computer Science, Procedia Economics and Finance
			\end{itemize}
		\end{block}
		\begin{block}{Strategy: Conference $\rightarrow$ Journal}
			\begin{itemize}
				\item Present at conference (get fast publication)
				\item Extend conference paper by 30-50\% new content
				\item Submit extended version to journal
				\item Two publications from one research effort
			\end{itemize}
		\end{block}
	\end{frame}
	
	% ============================================================
	% SLIDE 9: REVIEW PAPERS - HIGH CITATION POTENTIAL
	% ============================================================
	\begin{frame}{Review Papers - High Citation Potential}
		\scriptsize
		\begin{block}{What Are Review Papers?}
			\begin{itemize}
				\item Summarize and synthesize existing research on a topic
				\item Do not require new experimental data
				\item Cite 100-300 references
				\item Often invited by journal editors
			\end{itemize}
		\end{block}
		\begin{block}{Why Review Papers Get Many Citations}
			\begin{itemize}
				\item Anyone starting research in the field will cite the review
				\item Stay relevant for years (field overview)
				\item Often among most cited papers in a journal
				\item Can get 50-200+ citations within 2-3 years
			\end{itemize}
		\end{block}
		\begin{block}{How to Write a Review Paper as a Beginner}
			\begin{enumerate}
				\item Choose a trending topic (AI, LLMs, cybersecurity)
				\item Use Zotero to organize 100+ references
				\item Write systematic literature review
				\item Submit to mid-tier journal as "review article"
				\item Many journals accept unsolicited reviews
			\end{enumerate}
		\end{block}
	\end{frame}
	
	% ============================================================
	% SLIDE 10: BOOK CHAPTERS - ADDING DEPTH
	% ============================================================
	\begin{frame}{Book Chapters - Adding Depth to Portfolio}
		\scriptsize
		\begin{block}{What Are Book Chapters?}
			\begin{itemize}
				\item Contributions to edited volumes
				\item Published by academic presses (Springer, Elsevier, Wiley, IGI Global)
				\item Longer format (15-30 pages)
				\item Usually invited or proposed to editors
			\end{itemize}
		\end{block}
		\begin{block}{Why Book Chapters Help Your Profile}
			\begin{itemize}
				\item Count as a publication (different format than journals)
				\item Demonstrate ability to write at length
				\item Often indexed in Scopus and SpringerLink
				\item Can be cited just like journal papers
			\end{itemize}
		\end{block}
		\begin{block}{How to Get Book Chapter Invitations}
			\begin{itemize}
				\item Publish 2-3 journal papers first
				\item Editors will see your work and invite you
				\item Or propose to "Call for Chapters" announcements
				\item IGI Global and Springer frequently publish edited volumes
			\end{itemize}
		\end{block}
	\end{frame}
	
	% ============================================================
	% SLIDE 11: RESEARCH TOPICS THAT GET ATTENTION
	% ============================================================
	\begin{frame}{Research Topics That Get Attention}
		\scriptsize
		\begin{block}{Topics with High Government and Industry Interest}
			\centering
			\begin{tabular}{|p{5cm}|p{5cm}|}
				\hline
				\textbf{Topic Area} & \textbf{Specific Focus} \\
				\hline
				Artificial Intelligence & LLMs, GenAI safety, AI governance, algorithmic bias \\
				\hline
				Cybersecurity & Zero-trust, AI security, critical infrastructure \\
				\hline
				FinTech & Crypto regulation, digital assets, compliance \\
				\hline
				Workforce Development & AI retraining, digital literacy, upskilling \\
				\hline
				Financial Stability & Risk management, fraud detection, resilience \\
				\hline
				Healthcare AI & Clinical decision support, patient safety \\
				\hline
			\end{tabular}
		\end{block}
		\begin{block}{Where to Find Topic Ideas}
			\begin{itemize}
				\item Federal agency strategic plans (DHS, Treasury, Fed, SEC)
				\item Executive Orders (AI Executive Order 14110)
				\item Congressional hearings and reports
				\item Industry white papers and reports
			\end{itemize}
		\end{block}
	\end{frame}
	
	% ============================================================
	% SLIDE 12: TOPICS FROM FEDERAL AGENCY PRIORITIES
	% ============================================================
	\begin{frame}{Topics from Federal Agency Priorities}
		\scriptsize
		\begin{block}{Department of Treasury Priorities}
			\begin{itemize}
				\item Financial stability and risk modeling
				\item Cryptocurrency and digital asset regulation
				\item Sanctions compliance and enforcement
				\item Fraud detection and prevention
			\end{itemize}
		\end{block}
		\begin{block}{Federal Reserve Priorities}
			\begin{itemize}
				\item Payment systems modernization (FedNow)
				\item Climate risk and financial stability
				\item AI in banking supervision
				\item Community bank resilience
			\end{itemize}
		\end{block}
		\begin{block}{Department of Homeland Security Priorities}
			\begin{itemize}
				\item AI and cybersecurity workforce development
				\item Critical infrastructure protection
				\item Election security and disinformation
				\item Border technology and biometrics
			\end{itemize}
		\end{block}
		\begin{block}{Department of Education Priorities}
			\begin{itemize}
				\item AI literacy in K-12 and higher education
				\item Digital skills for workforce development
				\item Veteran retraining programs
				\item STEM education and diversity
			\end{itemize}
		\end{block}
	\end{frame}
	
	% ============================================================
	% SLIDE 13: CITATION STRATEGY - OVERVIEW
	% ============================================================
	\begin{frame}{Citation Strategy - Complete Overview}
		\scriptsize
		\begin{block}{Why Citations Matter}
			\begin{itemize}
				\item Citations prove your work is being used by others
				\item Show impact beyond yourself
				\item Demonstrate field recognition
				\item Are verifiable by third parties (Google Scholar, Scopus)
			\end{itemize}
		\end{block}
		\begin{block}{Types of Citations by Value}
			\centering
			\begin{tabular}{|p{4cm}|p{5.5cm}|}
				\hline
				\textbf{Citation Type} & \textbf{Value} \\
				\hline
				.gov citations (Federal Reserve, ERIC, Science.gov) & Highest \\
				\hline
				.edu citations (Harvard, MIT, Stanford) & Very High \\
				\hline
				Top journal citations (Nature, Science, IEEE Trans) & High \\
				\hline
				Scopus/WoS journal citations & Medium-High \\
				\hline
				Conference proceedings citations & Medium \\
				\hline
				Non-indexed journal citations & Low \\
				\hline
			\end{tabular}
		\end{block}
	\end{frame}
	
	% ============================================================
	% SLIDE 14: HOW TO INCREASE CITATION COUNT - BEFORE PUBLICATION
	% ============================================================
	\begin{frame}{Increase Citations - Before Publication}
		\scriptsize
		\begin{block}{Pre-Publication Strategies}
			\textbf{1. Choose High-Visibility Topics}
			\begin{itemize}
				\item AI, LLMs, cybersecurity, FinTech (trending now)
				\item Topics with active research communities
				\item Areas where government is investing
			\end{itemize}
			\textbf{2. Write Review Papers}
			\begin{itemize}
				\item Review papers get 2-3x more citations than original research
				\item Anyone new to field will cite your review
			\end{itemize}
			\textbf{3. Optimize Title and Abstract}
			\begin{itemize}
				\item Include keywords that people search for
				\item Make title specific and descriptive
				\item Write clear, searchable abstract
			\end{itemize}
			\textbf{4. Choose Open Access Journals}
			\begin{itemize}
				\item Anyone can read = more potential citations
				\item OA papers get 20-50\% more citations
			\end{itemize}
		\end{block}
	\end{frame}
	
	% ============================================================
	% SLIDE 15: HOW TO INCREASE CITATIONS - AFTER PUBLICATION
	% ============================================================
	\begin{frame}{Increase Citations - After Publication}
		\scriptsize
		\begin{block}{Post-Publication Strategies}
			\textbf{1. Upload to Preprint Servers}
			\begin{itemize}
				\item arXiv, SSRN, OSF, ResearchGate
				\item Citations start before journal publication
				\item Preprints get DOI and become citable
			\end{itemize}
			\textbf{2. Share on Social Media}
			\begin{itemize}
				\item LinkedIn: Share with professional network
				\item Twitter X: Use hashtags (AI Cybersecurity)
				\item ResearchGate: Upload full text
			\end{itemize}
			\textbf{3. Email Colleagues}
			\begin{itemize}
				\item Send PDF to researchers who cited similar papers
				\item Ask: "Would you be interested in this related work?"
				\item Offer to share data or code
			\end{itemize}
			\textbf{4. Present at Conferences}
			\begin{itemize}
				\item Conference presentations lead to citations
				\item Attendees download and cite your work
				\item Create slides with your paper citation
			\end{itemize}
		\end{block}
	\end{frame}
	
	% ============================================================
	% SLIDE 16: .GOV CITATIONS - THE GOLD STANDARD
	% ============================================================
	\begin{frame}{.gov Citations - The Gold Standard}
		\scriptsize
		\begin{block}{What Are .gov Citations?}
			\begin{itemize}
				\item Your work cited by U.S. federal government publications
				\item Examples: Federal Reserve, Bureau of Labor Statistics, ERIC, Science.gov
				\item .gov domain is verifiable by USCIS
			\end{itemize}
		\end{block}
		\begin{block}{Why .gov Citations Are Most Valuable}
			\begin{itemize}
				\item Direct evidence of national importance
				\item Federal agencies don't cite irrelevant work
				\item One .gov citation can be worth 50 regular citations
			\end{itemize}
		\end{block}
		\begin{block}{Where .gov Citations Appear}
			\begin{itemize}
				\item Federal Reserve Finance and Economics Discussion Series (FEDS)
				\item Bureau of Labor Statistics Monthly Labor Review
				\item ERIC (Department of Education) indexed papers
				\item Science.gov indexed research
				\item Congressional Research Service reports
				\item Government Accountability Office reports
			\end{itemize}
		\end{block}
	\end{frame}
	
	% ============================================================
	% SLIDE 17: FEDERAL RESERVE CITATIONS
	% ============================================================
	\begin{frame}{Federal Reserve Citations - Getting Cited by the Fed}
		\scriptsize
		\begin{block}{What is the FEDS Series?}
			\begin{itemize}
				\item Finance and Economics Discussion Series
				\item Official working paper series of Federal Reserve Board
				\item Peer-reviewed research on U.S. financial policy
			\end{itemize}
		\end{block}
		\begin{block}{How to Get Cited by Federal Reserve}
			\begin{enumerate}
				\item Publish on topics relevant to Fed: monetary policy, banking, financial stability, payments systems, AI in finance
				\item Make your paper available as preprint (SSRN, arXiv)
				\item Fed economists search these repositories
				\item If relevant, they will cite your work
			\end{enumerate}
			\begin{block}{Topics Fed Economists Cite}
				\begin{itemize}
					\item Financial risk modeling and stress testing
					\item Machine learning in banking supervision
					\item Payment systems and FedNow
					\item Climate risk and financial stability
					\item Cryptocurrency and digital assets
				\end{itemize}
			\end{block}
		\end{frame}
		
		% ============================================================
		% SLIDE 18: ERIC.GOV - DEPARTMENT OF EDUCATION INDEXING
		% ============================================================
		\begin{frame}{ERIC.gov - Department of Education Indexing}
			\scriptsize
			\begin{block}{What is ERIC?}
				\begin{itemize}
					\item Education Resources Information Center
					\item Digital library of U.S. Department of Education
					\item .gov domain authority
				\end{itemize}
			\end{block}
			\begin{block}{How to Get Indexed in ERIC}
				\textbf{Method 1: Journal Publication}
				\begin{itemize}
					\item Publish in ERIC-approved journals
					\item Journal automatically submits content
					\item Focus: education, workforce, training, digital literacy
				\end{itemize}
				\textbf{Method 2: Direct Submission}
				\begin{itemize}
					\item Submit unpublished research, technical reports, dissertations
					\item Must meet quality standards
					\item Submit via ERIC online portal
				\end{itemize}
			\end{block}
			\begin{block}{Best Topics for ERIC Indexing}
				\begin{itemize}
					\item AI literacy and digital skills training
					\item Workforce development and retraining
					\item Veteran education programs
					\item Online learning effectiveness
					\item STEM education
				\end{itemize}
			\end{block}
		\end{frame}
		
		% ============================================================
		% SLIDE 19: SCIENCE.GOV INDEXING
		% ============================================================
		\begin{frame}{Science.gov - Gateway to Federal Science}
			\scriptsize
			\begin{block}{What is Science.gov?}
				\begin{itemize}
					\item Official U.S. government portal for scientific information
					\item Searches 60+ federal databases simultaneously
					\item Managed by DOE, with participation from 14 federal agencies
				\end{itemize}
			\end{block}
			\begin{block}{How to Get Indexed in Science.gov}
				\textbf{Pathways to Science.gov:}
				\begin{itemize}
					\item PubMed (NIH) $\rightarrow$ Science.gov
					\item ERIC (Dept of Education) $\rightarrow$ Science.gov
					\item NASA ADS $\rightarrow$ Science.gov
					\item DOE OSTI $\rightarrow$ Science.gov
					\item USDA $\rightarrow$ Science.gov
				\end{itemize}
				\begin{block}{What This Means for You}
					\begin{itemize}
						\item Publish in journals indexed by PubMed, ERIC, or NASA ADS
						\item Your work automatically becomes discoverable through Science.gov
						\item Creates .gov discoverability without direct agency citation
					\end{itemize}
				\end{block}
			\end{frame}
			
			% ============================================================
			% SLIDE 20: BUREAU OF LABOR STATISTICS CITATIONS
			% ============================================================
			\begin{frame}{BLS Citations - Labor Market Impact}
				\scriptsize
				\begin{block}{What is BLS?}
					\begin{itemize}
						\item Bureau of Labor Statistics
						\item U.S. Department of Labor agency
						\item Publishes Monthly Labor Review, working papers, reports
					\end{itemize}
				\end{block}
				\begin{block}{How to Get Cited by BLS}
					\begin{enumerate}
						\item Publish on labor economics, workforce trends, employment statistics
						\item Submit comments to BLS RFIs
						\item Get referenced in Monthly Labor Review
						\item Collaborate with BLS researchers
					\end{enumerate}
					\begin{block}{Topics BLS Covers}
						\begin{itemize}
							\item Employment projections and trends
							\item Wage and compensation analysis
							\item Occupational outlook
							\item Workforce demographics
							\item Productivity and technology impact on jobs
						\end{itemize}
					\end{block}
				\end{frame}
				
				% ============================================================
				% SLIDE 21: .EDU CITATIONS - UNIVERSITY RECOGNITION
				% ============================================================
				\begin{frame}{.edu Citations - University Recognition}
					\scriptsize
					\begin{block}{What Are .edu Citations?}
						\begin{itemize}
							\item Your work cited by U.S. university researchers
							\item URL contains .edu domain (verifiable)
							\item Examples: Harvard, MIT, Stanford, Berkeley
						\end{itemize}
					\end{block}
					\begin{block}{Why .edu Citations Matter}
						\begin{itemize}
							\item Demonstrates academic recognition in the U.S.
							\item Verifiable by USCIS through .edu domain
							\item Proves work is used in U.S. education and research
						\end{itemize}
					\end{block}
					\begin{block}{How to Get .edu Citations}
						\begin{itemize}
							\item Publish in reputable journals (Scopus, WoS)
							\item Upload to institutional repositories (Harvard DASH, MIT DSpace)
							\item Present at university seminars
							\item Collaborate with US university researchers
							\item Create course-relevant materials that professors adopt
						\end{itemize}
					\end{block}
				\end{frame}
				
				% ============================================================
				% SLIDE 22: TRACKING .EDU CITATIONS
				% ============================================================
				\begin{frame}{Tracking .edu Citations}
					\scriptsize
					\begin{block}{Google Scholar Method}
						\begin{enumerate}
							\item Go to Google Scholar
							\item Click "My Profile"
							\item Click "Cited by" number for your paper
							\item Look at URLs of citing works
							\item Filter for URLs containing ".edu"
							\item Save these citations to your portfolio
						\end{enumerate}
					\end{block}
					\begin{block}{Advanced Search Method}
						\begin{verbatim}
							"your paper title" site:.edu
						\end{verbatim}
						\begin{itemize}
							\item Searches all .edu domains for your paper
							\item Reveals citations you may have missed
							\item Check every 3 months for new citations
						\end{itemize}
					\end{block}
					\begin{block}{Citations from Top US Universities}
						\begin{itemize}
							\item Harvard (.harvard.edu)
							\item MIT (.mit.edu)
							\item Stanford (.stanford.edu)
							\item Berkeley (.berkeley.edu)
							\item Caltech (.caltech.edu)
							\item University of Michigan (.umich.edu)
						\end{itemize}
					\end{block}
				\end{frame}
				
				% ============================================================
				% SLIDE 23: PEER REVIEW STRATEGY - OVERVIEW
				% ============================================================
				\begin{frame}{Peer Review Strategy - Complete Overview}
					\scriptsize
					\begin{block}{What is Peer Review?}
						\begin{itemize}
							\item Evaluating manuscripts submitted to journals
							\item Judging quality, originality, and significance
							\item Unpaid service to the academic community
						\end{itemize}
					\end{block}
					\begin{block}{Why Peer Review Builds Your Profile}
						\begin{itemize}
							\item Proves you are recognized as an expert
							\item Journal editors invite only qualified researchers
							\item Creates verifiable record (Web of Science, ORCID)
							\item Demonstrates "judging the work of others"
						\end{itemize}
					\end{block}
					\begin{block}{Two-Stage Path to Peer Review}
						\begin{enumerate}
							\item \textbf{Stage 1: Publish first} (2-3 papers in a journal)
							\item \textbf{Stage 2: Get invited to review} (editors notice active authors)
						\end{enumerate}
					\end{block}
				\end{frame}
				
				% ============================================================
				% SLIDE 24: HOW TO GET PEER REVIEW INVITATIONS
				% ============================================================
				\begin{frame}{How to Get Peer Review Invitations}
					\scriptsize
					\begin{block}{Method 1: Publish Actively}
						\begin{itemize}
							\item Publish 2-3 papers in a target journal
							\item Editors see your name repeatedly
							\item You become known as an active researcher in that field
							\item Editors will invite you to review
						\end{itemize}
					\end{block}
					\begin{block}{Method 2: Register on Journal Portals}
						\begin{itemize}
							\item Create accounts on ScholarOne (many journals use this)
							\item Register on Editorial Manager (Elsevier journals)
							\item Complete your reviewer profile with keywords
							\item Set yourself as "available for review"
						\end{itemize}
					\end{block}
					\begin{block}{Method 3: Email Journal Editors}
						\begin{itemize}
							\item Professional email: "I have published in your journal and would be happy to review"
							\item Attach your CV (2 pages max)
							\item List your areas of expertise
							\item Offer to review 2-3 papers per year
						\end{itemize}
					\end{block}
				\end{frame}
				
				% ============================================================
				% SLIDE 25: WEB OF SCIENCE REVIEWER RECOGNITION
				% ============================================================
				\begin{frame}{Web of Science Reviewer Recognition}
					\scriptsize
					\begin{block}{What is WoS Reviewer Recognition?}
						\begin{itemize}
							\item Formerly called Publons (now part of Web of Science)
							\item Centralized record of all your peer reviews
							\item Verified by journals (cannot be faked)
							\item Permanent record even if you change institutions
						\end{itemize}
					\end{block}
					\begin{block}{How to Set Up}
						\begin{enumerate}
							\item Go to webofscience.com
							\item Create free account
							\item Complete your profile (photo, affiliation, expertise)
							\item Add your publications (search by DOI)
							\item Link your ORCID
							\item Enable "Reviewer Recognition" to automatically track reviews
						\end{enumerate}
					\end{block}
					\begin{block}{What WoS Shows to USCIS}
						\begin{itemize}
							\item Number of reviews completed
							\item Names of journals you reviewed for
							\item Time period of reviews (sustained activity)
							\item Verification from journals
						\end{itemize}
					\end{block}
				\end{frame}
				
				% ============================================================
				% SLIDE 26: ORCID FOR PEER REVIEW VERIFICATION
				% ============================================================
				\begin{frame}{ORCID for Peer Review Verification}
					\scriptsize
					\begin{block}{What is ORCID?}
						\begin{itemize}
							\item Open Researcher and Contributor ID
							\item 16-digit permanent digital identifier
							\item Used by 1,000+ journals and publishers
							\item Free and permanent
						\end{itemize}
					\end{block}
					\begin{block}{How ORCID Tracks Peer Review}
						\begin{itemize}
							\item Journals can automatically push review data to ORCID
							\item Creates verified record (no manual entry possible to fake)
							\item Includes: journal name, date, type of review
							\item Some journals include editor thank-you notes
						\end{itemize}
					\end{block}
					\begin{block}{How to Connect ORCID to Journals}
						\begin{enumerate}
							\item Create ORCID profile (orcid.org)
							\item When submitting papers, add ORCID to your author profile
							\item When reviewing, provide ORCID to journal
							\item Journal pushes review credit to your ORCID record
						\end{enumerate}
					\end{block}
				\end{frame}
				
				% ============================================================
				% SLIDE 27: BUILDING A REVIEW PORTFOLIO - TIMELINE
				% ============================================================
				\begin{frame}{Building a Review Portfolio - Timeline}
					\scriptsize
					\begin{block}{Year 1: Foundation}
						\begin{itemize}
							\item Focus on publishing 2-3 papers
							\item Register on journal portals (ScholarOne, Editorial Manager)
							\item Complete reviewer profiles
							\item Target: Receive first 1-2 review invitations
						\end{itemize}
					\end{block}
					\begin{block}{Year 2: Building Volume}
						\begin{itemize}
							\item Complete 5-10 reviews across 3-4 journals
							\item Add ORCID and WoS reviewer recognition
							\item Target: 10-15 reviews total
						\end{itemize}
					\end{block}
					\begin{block}{Year 3+: Advanced Recognition}
						\begin{itemize}
							\item Complete 20-30+ reviews across 5-8 journals
							\item Apply for editorial board positions
							\item Guest edit special issues
							\item Mentor other reviewers
						\end{itemize}
					\end{block}
				\end{frame}
				
				% ============================================================
				% SLIDE 28: AWARD STRATEGY - BUILDING RECOGNITION
				% ============================================================
				\begin{frame}{Award Strategy - Building Recognition}
					\scriptsize
					\begin{block}{Types of Awards to Target}
						\begin{itemize}
							\item \textbf{Best Paper Awards:} At conferences you attend
							\item \textbf{Emerging Researcher Awards:} Many professional societies have these
							\item \textbf{Young Investigator Awards:} For researchers under 40
							\item \textbf{Fellowships:} Professional society fellows (IEEE, ACM, etc.)
							\item \textbf{Teaching/Service Awards:} From your university or employer
						\end{itemize}
					\end{block}
					\begin{block}{How to Get Awards}
						\begin{enumerate}
							\item Submit your best paper to conferences with best paper competitions
							\item Nominate yourself for society awards (many allow self-nomination)
							\item Ask mentors or collaborators to nominate you
							\item Build a strong publication record first (awards follow publications)
						\end{enumerate}
					\end{block}
					\begin{block}{Awards to Consider (Examples)}
						\begin{itemize}
							\item Globee Awards (business and technology)
							\item Titan Business Awards
							\item Stevie Awards
							\item IEEE Outstanding Young Professional Award
							\item ACM Doctoral Dissertation Award (if PhD completed)
						\end{itemize}
					\end{block}
				\end{frame}
				
				% ============================================================
				% SLIDE 29: AWARD DOCUMENTATION FOR PORTFOLIO
				% ============================================================
				\begin{frame}{Award Documentation for Your Portfolio}
					\scriptsize
					\begin{block}{What to Keep for Each Award}
						\begin{itemize}
							\item Award certificate (PDF)
							\item Announcement email or letter
							\item News article about the award
							\item Conference program showing award (if conference award)
							\item Photo of award ceremony (if available)
							\item List of criteria and selection process
						\end{itemize}
					\end{block}
					\begin{block}{Why Awards Build Your Profile}
						\begin{itemize}
							\item Third-party recognition of excellence
							\item Verifiable by USCIS (award website, news articles)
							\item Demonstrates your work stands out from peers
						\end{itemize}
					\end{block}
					\begin{block}{Where to Find Award Opportunities}
						\begin{itemize}
							\item Professional society websites (IEEE, ACM, INFORMS, ASA)
							\item Conference websites (best paper competitions)
							\item Industry award programs (Globee, Stevie, Titan)
							\item University research offices
						\end{itemize}
					\end{block}
				\end{frame}
				
				% ============================================================
				% SLIDE 30: PREPRINT SERVERS - ACCELERATE CITATIONS
				% ============================================================
				\begin{frame}{Preprint Servers - Accelerate Your Citations}
					\scriptsize
					\begin{block}{What Are Preprint Servers?}
						\begin{itemize}
							\item Online platforms to share research before journal publication
							\item Assign permanent DOI to your work
							\item Work becomes citable immediately
							\item Free to upload and access
						\end{itemize}
					\end{block}
					\begin{block}{Major Preprint Servers}
						\centering
						\begin{tabular}{|p{4cm}|p{5.5cm}|}
							\hline
							\textbf{Server} & \textbf{Best For} \\
							\hline
							arXiv.org & Physics, math, computer science, quantitative finance \\
							\hline
							SSRN (Elsevier) & Social sciences, economics, law, finance \\
							\hline
							OSF Preprints & Interdisciplinary research \\
							\hline
							ResearchGate & All fields (upload as preprint) \\
							\hline
							medRxiv/bioRxiv & Health and life sciences \\
							\hline
						\end{tabular}
					\end{block}
					\begin{block}{Why Preprints Build Your Profile}
						\begin{itemize}
							\item Citations start immediately (not after 6-12 months)
							\item Google Scholar indexes preprints
							\item Your work reaches audience faster
							\item Creates timestamped record of your research
						\end{itemize}
					\end{block}
				\end{frame}
				
				% ============================================================
				% SLIDE 31: SSRN FOR FINANCE AND ECONOMICS RESEARCH
				% ============================================================
				\begin{frame}{SSRN - For Finance, Economics, and Law Research}
					\scriptsize
					\begin{block}{What is SSRN?}
						\begin{itemize}
							\item Social Science Research Network
							\item Owned by Elsevier
							\item Leading preprint server for economics, finance, law
							\item Used by Federal Reserve, Treasury, SEC researchers
						\end{itemize}
					\end{block}
					\begin{block}{How SSRN Builds Your Profile}
						\begin{itemize}
							\item Your paper is searchable by government economists
							\item Download counts show interest in your work
							\item SSRN rankings (top 10\% downloads) demonstrate impact
							\item Many .gov citations originate from SSRN discovery
						\end{itemize}
					\end{block}
					\begin{block}{How to Upload to SSRN}
						\begin{enumerate}
							\item Create free account (ssrn.com)
							\item Click "Submit a Paper"
							\item Upload PDF and metadata
							\item Choose subject area (e.g., FinTech, AI in Finance)
							\item Paper appears within 24-48 hours
						\end{enumerate}
					\end{block}
				\end{frame}
				
				% ============================================================
				% SLIDE 32: OSF PREPRINTS - PERMANENT DOI
				% ============================================================
				\begin{frame}{OSF Preprints - Permanent DOIs for Your Work}
					\scriptsize
					\begin{block}{What is OSF?}
						\begin{itemize}
							\item Open Science Framework (osf.io)
							\item Maintained by Center for Open Science
							\item Receives funding from NIH and NSF
						\end{itemize}
					\end{block}
					\begin{block}{Why OSF is Strategic for Profile Building}
						\begin{itemize}
							\item Every preprint gets permanent DOI (citable immediately)
							\item OSF preprints are harvested by SHARE (U.S. Department of Education project)
							\item Your work becomes discoverable through .gov portals
							\item Free and no waiting for journal acceptance
						\end{itemize}
					\end{block}
					\begin{block}{How to Use OSF}
						\begin{enumerate}
							\item Create free account at osf.io
							\item Click "Create a New Project"
							\item Upload your preprint PDF
							\item OSF assigns DOI automatically
							\item Share DOI on LinkedIn, email signature, Google Scholar
						\end{enumerate}
					\end{block}
				\end{frame}
				
				% ============================================================
				% SLIDE 33: EBSCO AND PROQUEST INDEXING
				% ============================================================
				\begin{frame}{EBSCO and ProQuest - Commercial Database Indexing}
					\scriptsize
					\begin{block}{What is EBSCO?}
						\begin{itemize}
							\item Commercial database used by 1,000s of university libraries
							\item U.S. Commerce Research Library uses EBSCO Business Source Complete
							\item Papers in EBSCO are searchable by US government policymakers
						\end{itemize}
					\end{block}
					\begin{block}{How to Get EBSCO Indexed}
						\begin{itemize}
							\item Publish in journals indexed by EBSCO
							\item Check: https://about.ebsco.com/title-lists
							\item Focus on business, economics, technology journals
						\end{itemize}
					\end{block}
					\begin{block}{What is ProQuest?}
						\begin{itemize}
							\item Dissertations, theses, conference proceedings
							\item Submit through your university library
							\item Creates permanent, citable record of your work
						\end{itemize}
					\end{block}
				\end{frame}
				
				% ============================================================
				% SLIDE 34: U.S. COMMERCE LIBRARY INDEXING
				% ============================================================
				\begin{frame}{U.S. Commerce Library Indexing - .gov Discovery}
					\scriptsize
					\begin{block}{What is the Commerce Research Library?}
						\begin{itemize}
							\item Official library of the U.S. Department of Commerce
							\item Indexes journals through EBSCO Business Source Complete
							\item Your work becomes searchable by Commerce Department policymakers
						\end{itemize}
					\end{block}
					\begin{block}{How to Get Indexed in Commerce Library}
						\begin{enumerate}
							\item Publish in EBSCO Business Source Complete journals
							\item Ensure paper focuses on: trade, finance, economics, technology, business
							\item Wait 3-6 months for indexing
							\item Verify at search.library.doc.gov
						\end{enumerate}
					\end{block}
					\begin{block}{What to Document for Your Portfolio}
						\begin{itemize}
							\item Screenshot of search results showing your paper
							\item URL (starts with .gov)
							\item Date of search
						\end{itemize}
					\end{block}
				\end{frame}
				
				% ============================================================
				% SLIDE 35: GOOGLE SCHOLAR PROFILE OPTIMIZATION
				% ============================================================
				\begin{frame}{Google Scholar Profile - Complete Optimization}
					\scriptsize
					\begin{block}{Why Google Scholar Matters}
						\begin{itemize}
							\item USCIS officers check Google Scholar
							\item Publicly verifiable citation counts
							\item Shows citation growth over time
							\item Free and easy to maintain
						\end{itemize}
					\end{block}
					\begin{block}{How to Optimize Your Profile}
						\begin{enumerate}
							\item Add professional photo (optional but recommended)
							\item Use consistent name across all publications
							\item Add affiliation (.edu if you have one)
							\item Add research keywords (helps discovery)
							\item Make profile PUBLIC (USCIS must be able to view)
							\item Enable email updates (notified when cited)
						\end{enumerate}
					\end{block}
					\begin{block}{What to Include in Your Profile}
						\begin{itemize}
							\item All journal papers (including preprints)
							\item Conference proceedings
							\item Book chapters
							\item Technical reports
							\item Patents (if any)
						\end{itemize}
					\end{block}
				\end{frame}
				
				% ============================================================
				% SLIDE 36: RESEARCHGATE FOR NETWORKING
				% ============================================================
				\begin{frame}{ResearchGate - Professional Network for Researchers}
					\scriptsize
					\begin{block}{What is ResearchGate?}
						\begin{itemize}
							\item Professional network for scientists and researchers
							\item 20+ million users worldwide
							\item Share papers, ask questions, follow research
						\end{itemize}
					\end{block}
					\begin{block}{How ResearchGate Builds Your Profile}
						\begin{itemize}
							\item Upload full text of your papers (check copyright)
							\item See who is reading your work (download stats)
							\item Receive recommendations from other researchers
							\item RG Score (research impact metric)
							\item Global visibility for collaborations
						\end{itemize}
					\end{block}
					\begin{block}{What to Track on ResearchGate}
						\begin{itemize}
							\item Number of reads (indicates interest)
							\item Citation count (verified)
							\item Recommendations from peers
							\item International collaboration network
							\item Screenshot profile for your portfolio
						\end{itemize}
					\end{block}
				\end{frame}
				
				% ============================================================
				% SLIDE 37: LINKEDIN FOR RESEARCH VISIBILITY
				% ============================================================
				\begin{frame}{LinkedIn - Professional Visibility}
					\scriptsize
					\begin{block}{Why LinkedIn Matters for Researchers}
						\begin{itemize}
							\item Recruiters and collaborators search LinkedIn
							\item Showcase publications with links and PDFs
							\item Build network with other researchers
						\end{itemize}
					\end{block}
					\begin{block}{How to Optimize Your LinkedIn Profile}
						\begin{enumerate}
							\item Add "Publications" section with all papers
							\item Include DOI links (verifiable)
							\item Upload PDFs when possible
							\item Add "Projects" section for research initiatives
							\item Request recommendations from collaborators
						\end{enumerate}
					\end{block}
					\begin{block}{Sharing Research on LinkedIn}
						\begin{itemize}
							\item Post when new paper is published
							\item Share preprints with a brief summary
							\item Tag co-authors and collaborators
							\item Use relevant hashtags (#AI #Research)
							\item Engage with comments (increases visibility)
						\end{itemize}
					\end{block}
				\end{frame}
				
				% ============================================================
				% SLIDE 38: MEDIA AND OUTREACH STRATEGY
				% ============================================================
				\begin{frame}{Media and Outreach - Building Public Recognition}
					\scriptsize
					\begin{block}{Why Media Coverage Helps Your Profile}
						\begin{itemize}
							\item Demonstrates your work has public interest
							\item Third-party validation of importance
							\item Creates verifiable evidence (news articles, podcasts)
						\end{itemize}
					\end{block}
					\begin{block}{How to Get Media Coverage}
						\begin{enumerate}
							\item Write a press release for your paper
							\item Pitch to university communications office (if affiliated)
							\item Contact science journalists who cover your topic
							\item Offer to write guest blog posts
							\item Start your own blog or newsletter
						\end{enumerate}
					\end{block}
					\begin{block}{Low-Effort Media Opportunities}
						\begin{itemize}
							\item LinkedIn articles (publish directly on LinkedIn)
							\item Medium.com (free blogging platform)
							\item Podcast guest appearances (many podcasts seek experts)
							\item Webinar presentations (record and share)
							\item Twitter/X threads summarizing your paper
						\end{itemize}
					\end{block}
				\end{frame}
				
				% ============================================================
				% SLIDE 39: CREATING TEACHING/TRAINING MATERIALS
				% ============================================================
				\begin{frame}{Creating Teaching and Training Materials}
					\scriptsize
					\begin{block}{Why Teaching Materials Build Your Profile}
						\begin{itemize}
							\item Demonstrates ability to communicate complex topics
							\item Creates additional citations (when others use your materials)
							\item Shows workforce development contribution
						\end{itemize}
					\end{block}
					\begin{block}{Types of Teaching/Training Materials}
						\begin{itemize}
							\item Online courses (Udemy, Coursera, YouTube)
							\item Workshop slide decks (share on SlideShare)
							\item Tutorial videos (YouTube)
							\item Open educational resources (OER)
							\item Guest lectures at universities
						\end{itemize}
					\end{block}
					\begin{block}{How Training Materials Get Cited}
						\begin{itemize}
							\item Other instructors link to your materials
							\item Students cite your course in their work
							\item Universities include your resources in syllabi
							\item Creates .edu citations when used in courses
						\end{itemize}
					\end{block}
				\end{frame}
				
				% ============================================================
				% SLIDE 40: UDEMY COURSES FOR WORKFORCE IMPACT
				% ============================================================
				\begin{frame}{Udemy Courses - Demonstrating Workforce Impact}
					\scriptsize
					\begin{block}{Why Udemy Helps Your Profile}
						\begin{itemize}
							\item Quantifiable metrics (enrollments, completions, ratings)
							\item Demonstrates public reach and impact
							\item Verifiable by third parties
						\end{itemize}
					\end{block}
					\begin{block}{What to Document from Your Course}
						\begin{itemize}
							\item Number of students enrolled (total and current)
							\item Geographic distribution (nationwide or global)
							\item Completion rates
							\item Student reviews and ratings
							\item Certificates issued
						\end{itemize}
					\end{block}
					\begin{block}{Best Topics for Udemy Courses}
						\begin{itemize}
							\item AI literacy for professionals
							\item Cybersecurity fundamentals
							\item Data science for finance
							\item Digital skills for workforce development
							\item Programming languages (Python, R, SQL)
						\end{itemize}
					\end{block}
				\end{frame}
				
				% ============================================================
				% SLIDE 41: BUILDING GITHUB PORTFOLIO
				% ============================================================
				\begin{frame}{GitHub Portfolio - Code and Research Transparency}
					\scriptsize
					\begin{block}{Why GitHub Matters for Researchers}
						\begin{itemize}
							\item Showcases code underlying your research
							\item Demonstrates transparency and reproducibility
							\item Creates additional citable artifacts (with DOIs via Zenodo)
						\end{itemize}
					\end{block}
					\begin{block}{What to Put on GitHub}
						\begin{itemize}
							\item Code for data analysis and experiments
							\item LaTeX source for papers
							\item Datasets (if shareable)
							\item Tutorials and examples
							\item README files explaining your research
						\end{itemize}
					\end{block}
					\begin{block}{GitHub Metrics to Track}
						\begin{itemize}
							\item Stars (bookmarks from other users)
							\item Forks (others using your code)
							\item Downloads/releases
							\item Contributors (collaborators)
						\end{itemize}
					\end{block}
				\end{frame}
				
				% ============================================================
				% SLIDE 42: QUANTIFYING YOUR PROGRESS
				% ============================================================
				\begin{frame}{Quantifying Your Progress - Metrics Dashboard}
					\scriptsize
					\begin{block}{Publications Metrics to Track}
						\begin{itemize}
							\item Total number of publications
							\item Number by type (journal, conference, book chapter)
							\item Journal tiers (Level 0, Mid, High)
							\item Scopus/WoS indexed papers count
						\end{itemize}
					\end{block}
					\begin{block}{Citation Metrics to Track}
						\begin{itemize}
							\item Total Google Scholar citations
							\item Citation growth rate (quarter over quarter)
							\item h-index and i10-index
							\item Number of .gov citations
							\item Number of .edu citations
						\end{itemize}
					\end{block}
					\begin{block}{Review Metrics to Track}
						\begin{itemize}
							\item Total peer reviews completed
							\item Number of distinct journals reviewed for
							\item WoS verified reviews
							\item Editorial board positions
						\end{itemize}
					\end{block}
				\end{frame}
				
				% ============================================================
				% SLIDE 43: SAMPLE 12-MONTH PROGRESSION
				% ============================================================
				\begin{frame}{Sample 12-Month Progression - Starting from Zero}
					\scriptsize
					\begin{block}{Month 1-3: Foundation}
						\begin{itemize}
							\item Literature review on chosen topic
							\item Set up Google Scholar, ORCID, ResearchGate
							\item Write first review paper
							\item Submit to Level 0 or mid-tier journal
						\end{itemize}
					\end{block}
					\begin{block}{Month 4-6: First Publications}
						\begin{itemize}
							\item First paper accepted (1 publication)
							\item Submit second paper (original research)
							\item Register on journal portals for peer review
						\end{itemize}
					\end{block}
					\begin{block}{Month 7-9: Building Momentum}
						\begin{itemize}
							\item Second paper accepted (2 publications)
							\item Submit third paper (conference or journal)
							\item Receive first peer review invitation (review 1)
						\end{itemize}
					\end{block}
					\begin{block}{Month 10-12: Impact Emerging}
						\begin{itemize}
							\item Third paper accepted (3 publications)
							\item Complete 2-3 peer reviews
							\item First citations appear on Google Scholar
							\item Upload all papers as preprints
						\end{itemize}
					\end{block}
				\end{frame}
				
				% ============================================================
				% SLIDE 44: SAMPLE 24-MONTH PROGRESSION
				% ============================================================
				\begin{frame}{Sample 24-Month Progression - Advanced Profile}
					\scriptsize
					\begin{block}{Year 2 Targets}
						\begin{itemize}
							\item 8-12 total publications
							\item 5-8 Scopus indexed papers
							\item 30-100+ Google Scholar citations
							\item 10-20 peer reviews completed
							\item 1-2 .gov citations (if topics aligned with federal priorities)
							\item 3-5 .edu citations
							\item 1-2 awards or media mentions
						\end{itemize}
					\end{block}
					\begin{block}{Month 13-18: Scaling Up}
						\begin{itemize}
							\item Submit 2-3 more papers
							\item Present at 1-2 conferences
							\item Build citation network (papers being cited)
							\item Apply for editorial board positions
						\end{itemize}
					\end{block}
					\begin{block}{Month 19-24: Portfolio Complete}
						\begin{itemize}
							\item All publications complete
							\item Steady citation growth
							\item Peer review portfolio documented
							\item Expert letters from collaborators
							\item Complete research profile ready
						\end{itemize}
					\end{block}
				\end{frame}
				
				% ============================================================
				% SLIDE 45: COMMON MISTAKES TO AVOID
				% ============================================================
				\begin{frame}{Common Mistakes That Hurt Your Profile}
					\scriptsize
					\begin{block}{Publication Mistakes}
						\begin{itemize}
							\item \textcolor{red}{Publishing in predatory journals} (hurts credibility)
							\item \textcolor{red}{Paying high APCs (\$2,000+) for low-tier journals}
							\item \textcolor{red}{Not verifying Scopus indexing before submitting}
							\item \textcolor{red}{Only publishing in one journal (lack of breadth)}
						\end{itemize}
					\end{block}
					\begin{block}{Citation Mistakes}
						\begin{itemize}
							\item \textcolor{red}{Self-citing excessively} (looks manipulative)
							\item \textcolor{red}{Not tracking .gov and .edu citations}
							\item \textcolor{red}{Not uploading to preprint servers}
						\end{itemize}
					\end{block}
					\begin{block}{Profile Management Mistakes}
						\begin{itemize}
							\item \textcolor{red}{Private Google Scholar profile (USCIS can't see)}
							\item \textcolor{red}{No ORCID profile}
							\item \textcolor{red}{Not documenting progress with screenshots}
							\item \textcolor{red}{Inconsistent name across publications}
						\end{itemize}
					\end{block}
				\end{frame}
				
				% ============================================================
				% SLIDE 46: DOCUMENTATION BEST PRACTICES
				% ============================================================
				\begin{frame}{Documentation Best Practices}
					\scriptsize
					\begin{block}{What to Document and Save}
						\begin{itemize}
							\item [$\square$] PDF of every publication (final version)
							\item [$\square$] Acceptance emails from journals
							\item [$\square$] Screenshots of Google Scholar profile (quarterly)
							\item [$\square$] Screenshots of citations with .gov and .edu URLs
							\item [$\square$] Peer review invitation emails
							\item [$\square$] Peer review completion confirmations
							\item [$\square$] Web of Science reviewer recognition screenshots
							\item [$\square$] ORCID record screenshots
							\item [$\square$] Award certificates and announcements
							\item [$\square$] Media mentions (full articles, screenshots)
							\item [$\square$] Course enrollment data (Udemy, etc.)
						\end{itemize}
					\end{block}
				\end{frame}
				
				% ============================================================
				% SLIDE 47: TOOLS SUMMARY
				% ============================================================
				\begin{frame}{Tools Summary for Profile Building}
					\scriptsize
					\begin{block}{Essential Free Tools}
						\centering
						\begin{tabular}{|p{4cm}|p{5.5cm}|}
							\hline
							\textbf{Tool} & \textbf{Purpose} \\
							\hline
							Google Scholar & Citation tracking, profile visibility \\
							\hline
							ORCID & Permanent researcher ID \\
							\hline
							ResearchGate & Networking, paper sharing \\
							\hline
							Zotero & Reference management \\
							\hline
							LaTeX (Overleaf/TeXstudio) & Paper writing \\
							\hline
							Git/GitHub & Version control, code sharing \\
							\hline
							OSF Preprints & Free preprints with DOI \\
							\hline
							SSRN & Preprints for finance/economics \\
							\hline
						\end{tabular}
					\end{block}
					\begin{block}{Optional Paid Tools}
						\begin{itemize}
							\item Grammarly (editing, \$12/month)
							\item Overleaf Pro (collaboration, \$15/month)
							\item ChatGPT Plus (research drafting, \$20/month)
						\end{itemize}
					\end{block}
				\end{frame}
				
				% ============================================================
				% SLIDE 48: QUICK REFERENCE - PUBLICATION TARGETS
				% ============================================================
				\begin{frame}{Quick Reference - Publication Targets by Timeline}
					\scriptsize
					\begin{block}{6 Months (Minimum to Show Activity)}
						\begin{itemize}
							\item 1-2 publications (Level 0 or conference)
							\item 0-5 citations
							\item 0-1 peer reviews
						\end{itemize}
					\end{block}
					\begin{block}{12 Months (Competitive NIW Profile)}
						\begin{itemize}
							\item 4-6 publications (mix of Level 0, mid-tier, conference)
							\item 10-30 citations
							\item 2-5 peer reviews
							\item Google Scholar profile established
						\end{itemize}
					\end{block}
					\begin{block}{18-24 Months (Strong NIW / EB1A Starting Point)}
						\begin{itemize}
							\item 8-12 publications
							\item 50-150+ citations
							\item 10-20+ peer reviews
							\item 1-3 .gov or .edu citations
							\item ORCID and WoS reviewer profiles complete
						\end{itemize}
					\end{block}
				\end{frame}
				
				% ============================================================
				% SLIDE 49: FINAL CHECKLIST - PROFILE READY
				% ============================================================
				\begin{frame}{Final Checklist - Is Your Research Profile Ready?}
					\scriptsize
					\begin{block}{Publications Checklist}
						\begin{itemize}
							\item [$\square$] 5+ publications (minimum)
							\item [$\square$] At least 3 Scopus/WoS indexed papers
							\item [$\square$] Mix of journal and conference papers
							\item [$\square$] At least 1 review paper (if possible)
						\end{itemize}
					\end{block}
					\begin{block}{Citations Checklist}
						\begin{itemize}
							\item [$\square$] 30+ Google Scholar citations (minimum)
							\item [$\square$] Citation growth over 6+ months
							\item [$\square$] Screenshots of .gov or .edu citations (if any)
						\end{itemize}
					\end{block}
					\begin{block}{Peer Review Checklist}
						\begin{itemize}
							\item [$\square$] 5+ completed peer reviews
							\item [$\square$] Web of Science or ORCID record of reviews
							\item [$\square$] Reviews across 2+ distinct journals
						\end{itemize}
					\end{block}
					\begin{block}{Profile Management Checklist}
						\begin{itemize}
							\item [$\square$] Public Google Scholar profile
							\item [$\square$] Complete ORCID profile
							\item [$\square$] ResearchGate profile with papers
							\item [$\square$] Publications uploaded to preprint servers
						\end{itemize}
					\end{block}
				\end{frame}
				
				% ============================================================
				% SLIDE 50: CONCLUSION AND NEXT STEPS
				% ============================================================
				\begin{frame}{Conclusion and Next Steps}
					\scriptsize
					\begin{block}{The Complete Profile Building Framework}
						\begin{enumerate}
							\item \textbf{Publish strategically} - Level 0 $\rightarrow$ Mid-Tier $\rightarrow$ High-Tier
							\item \textbf{Maximize citations} - Preprints, open access, social sharing
							\item \textbf{Build review portfolio} - Publish first, then review
							\item \textbf{Target .gov and .edu} - Align research with federal priorities
							\item \textbf{Document everything} - Screenshots, URLs, timestamps
							\item \textbf{Use free tools} - Google Scholar, ORCID, OSF, ResearchGate
							\item \textbf{Track progress monthly} - Update your metrics dashboard
						\end{enumerate}
					\end{block}
					\begin{block}{Start Today - First 7 Days}
						\begin{enumerate}
							\item Day 1: Create Google Scholar profile
							\item Day 2: Create ORCID profile
							\item Day 3: Create ResearchGate account
							\item Day 4: Install Zotero
							\item Day 5: Identify 3 target journals
							\item Day 6: Write 500 words of literature review
							\item Day 7: Set up OSF account for preprints
						\end{enumerate}
					\end{block}
					\begin{block}{Final Thought}
						\textcolor{blue}{Consistency beats intensity. Thirty minutes every day builds a world-class research profile. Start today.}
					\end{block}
				\end{frame}
				
			\end{document}